\chapter{Colore dei sali}

Una molecola assorbe radiazione elettromagnetica quando l'energia del fotone coincide con la differenza fra due suoi livelli energetici.
A seconda della scala di energia in gioco si eccitano gradi di libertà diversi: le microonde muovono i livelli rotazionali, l'infrarosso quelli vibrazionali, mentre l'ultravioletto e il visibile promuovono gli elettroni da un orbitale all'altro.
Sono queste ultime, le transizioni elettroniche\mkeyword{transizione elettronica}, a determinare il colore.

La transizione meno costosa, e quindi la più rilevante, è quella che porta un elettrone dall'orbitale occupato di energia più alta, detto HOMO\mkeyword{HOMO}, al primo orbitale vuoto, detto LUMO\mkeyword{LUMO}, e richiede un fotone di energia
\begin{equation}
    \Delta E = h\nu = \frac{hc}{\lambda} .
    \label{eq:homo_lumo}
\end{equation}
Perché la sostanza appaia colorata occorre che questa transizione cada nell'intervallo visibile, all'incirca fra $380$ e $700 \ nm$, che per la \eqref{eq:homo_lumo} corrisponde a energie comprese fra $1{,}8$ e $3{,}3 \ eV$, cioè fra $170$ e $315 \ kJ/mol$.
È un intervallo piuttosto stretto, e ben più piccolo del salto fra un orbitale di legame e il corrispondente di antilegame: per questo la stragrande maggioranza dei composti è incolore, e il colore richiede un meccanismo che produca livelli molto ravvicinati.
Va aggiunto, per evitare un equivoco frequente, che il colore percepito non è quello assorbito ma il suo complementare: una soluzione che assorbe nel rosso appare azzurra, come accade allo ione rameico idratato.

\section{Ioni complessi e campo cristallino}

Il meccanismo che genera livelli sufficientemente ravvicinati è offerto dai metalli di transizione, i cui orbitali $d$ sono degeneri nell'atomo isolato e cessano di esserlo quando lo ione è circondato da altre specie.
Il caso tipico è quello degli ioni complessi\mkeyword{ione complesso}, nei quali uno ione metallico lega intorno a sé un certo numero di molecole o ioni, detti leganti\mkeyword{legante}, mediante legami di coordinazione; lo stesso avviene, in forma meno evidente, per uno ione inserito nel reticolo cristallino di un sale, dove il ruolo dei leganti è svolto dagli ioni circostanti.
Uno ione metallico non è mai realmente isolato, ed è questo il punto di partenza.

Se il campo elettrico prodotto dall'intorno avesse simmetria sferica, tutti e cinque gli orbitali $d$ sentirebbero la stessa repulsione e resterebbero degeneri, semplicemente innalzati in energia.
Il campo reale però non è sferico, ed è la sua simmetria a rimuovere la degenerazione\mkeyword{rimozione della degenerazione}.
In un campo ottaedrico i sei leganti si dispongono lungo gli assi cartesiani: gli orbitali $d_{z^2}$ e $d_{x^2-y^2}$, che puntano verso di essi, subiscono una repulsione maggiore e vengono destabilizzati, mentre $d_{xy}$, $d_{xz}$ e $d_{yz}$, che puntano fra i leganti, vengono stabilizzati.
La separazione fra i due gruppi si indica con $\Delta_o$, e poiché il baricentro energetico deve conservarsi i due orbitali alti salgono di $+0{,}6\,\Delta_o$ e i tre bassi scendono di $-0{,}4\,\Delta_o$, con bilancio complessivo
\begin{equation}
    2 \cdot (+0{,}6\,\Delta_o) + 3 \cdot (-0{,}4\,\Delta_o) = 0 .
    \label{eq:baricentro}
\end{equation}
Il punto essenziale è che $\Delta_o$ risulta molto minore della separazione fra orbitali di legame e di antilegame, e cade proprio nell'intervallo visibile: la transizione $d$--$d$ è dunque la transizione HOMO--LUMO dei complessi dei metalli di transizione, ed è l'origine del loro colore.

Il campo tetraedrico produce lo schema opposto, perché i quattro leganti non puntano lungo gli assi: si stabilizzano i due orbitali che nel caso ottaedrico erano alti e si destabilizzano gli altri tre, con una separazione minore, pari a $\Delta_t = \tfrac{4}{9}\Delta_o$, dal momento che i leganti sono meno numerosi e peggio orientati.
Il campo quadrato planare, più raro, si ottiene idealmente allontanando i due leganti lungo $z$ da un ottaedro: restano i quattro leganti nel piano $xy$, cosicché $d_{xz}$, $d_{yz}$ e $d_{xy}$ si stabilizzano, mentre $d_{x^2-y^2}$, che punta direttamente verso i leganti, viene fortemente destabilizzato e $d_{z^2}$ resta a energia intermedia, destabilizzato soltanto dal toro che la sua funzione d'onda possiede nel piano equatoriale.

\section{Configurazioni elettroniche e regola dei sali incolori}

Per applicare lo schema occorre conoscere la configurazione dello ione.
Nei cationi dei metalli di transizione gli orbitali $4s$ si trovano al di sopra dei $3d$, e sono quindi i primi a essere svuotati: uno ione bipositivo perde i due elettroni $s$ e conserva tutti i suoi elettroni $d$.
L'ordine di riempimento che si usa per gli atomi neutri, nel quale il $4s$ precede il $3d$, non vale dunque per la ionizzazione, e la contraddizione è solo apparente, perché l'ordine dei livelli dipende dalla carica dello ione.

Ne segue una regola semplice.
Soltanto le configurazioni comprese fra $d^1$ e $d^9$ possiedono orbitali $d$ parzialmente occupati e possono dare transizioni $d$--$d$; le configurazioni $d^0$ e $d^{10}$ non lo possono, perché nel primo caso non c'è nessun elettrone da promuovere e nel secondo non c'è nessun posto libero dove promuoverlo.
I cationi dei gruppi principali sono tutti $d^0$ oppure $d^{10}$, ed è per questo che i loro sali sono incolori; lo stesso vale per il cloruro rameoso, dove il rame $\mathrm{Cu^+}$ è $d^{10}$ ed è proprio la simmetria chiusa di quel guscio a renderlo stabile e il sale bianco.

\section{Alto spin, basso spin e serie spettrochimica}

Quando gli elettroni $d$ sono più di tre nasce una competizione fra due costi: promuovere un elettrone al livello superiore costa $\Delta$, mentre appaiarlo in un orbitale già occupato costa l'energia di appaiamento $P$, dovuta alla repulsione fra i due elettroni nella stessa regione di spazio.
Il caso di scuola è lo ione $\mathrm{Fe^{3+}}$, che è $d^5$: se $\Delta < P$ i cinque elettroni occupano singolarmente tutti e cinque gli orbitali, e si ha una configurazione ad alto spin\mkeyword{alto spin}; se invece $\Delta > P$ conviene appaiarli nei tre orbitali bassi, e si ha una configurazione a basso spin\mkeyword{basso spin}.
La scelta dipende dai leganti, poiché sono essi a fissare $\Delta$: leganti che generano un campo intenso favoriscono il basso spin.

L'ordine sperimentale dei leganti secondo il campo che generano prende il nome di serie spettrochimica\mkeyword{serie spettrochimica},
\[
    \mathrm{I^-} < \mathrm{Cl^-} < \mathrm{F^-} < \mathrm{OH^-} < \mathrm{H_2O} < \mathrm{NH_3} < \mathrm{NO_2^-} < \mathrm{CN^-} < \mathrm{CO} ,
\]
e ha una conseguenza immediatamente osservabile: cambiando il legante attorno allo stesso ione metallico cambia $\Delta$, e con esso la lunghezza d'onda assorbita e il colore del composto.
Un campo più forte assorbe a lunghezze d'onda minori.

Su questo punto conviene però essere onesti riguardo ai limiti del modello.
La teoria del campo cristallino è puramente elettrostatica, e come tale prevederebbe che i leganti più carichi e più piccoli generino i campi più intensi; la serie spettrochimica dice invece il contrario, ponendo lo ioduro in fondo e molecole neutre come $\mathrm{CO}$ in cima.
La spiegazione richiede di abbandonare l'elettrostatica e di ragionare in termini di orbitali molecolari: leganti come $\mathrm{CO}$ e $\mathrm{CN^-}$ possiedono orbitali $\pi$ di antilegame vuoti e di energia opportuna, che interagendo con gli orbitali $d$ occupati del metallo li abbassano ulteriormente, allargando $\Delta$.
Il campo cristallino resta un modello utile per contare i livelli e prevedere le degenerazioni, ma non spiega la loro spaziatura.

Una seconda precisazione riguarda l'intensità del colore.
Le transizioni $d$--$d$ sono formalmente proibite dalle regole di selezione, poiché non cambiano la parità dell'orbitale, e diventano possibili soltanto grazie alle vibrazioni che deformano momentaneamente la simmetria: per questo i complessi dei metalli di transizione sono tipicamente colorati ma di colore tenue.
I colori intensissimi, come il viola del permanganato, hanno origine diversa e sono dovuti a transizioni di trasferimento di carica dal legante al metallo.
La verifica sperimentale è netta: $\mathrm{MnO_4^-}$ contiene manganese al numero di ossidazione $+7$, cioè $d^0$, e secondo il ragionamento precedente dovrebbe essere incolore, mentre è una delle sostanze più intensamente colorate che si conoscano.
